\begin{frame}{2.2절 정리}
  \begin{block}{핵심}
    \kb{낙관적 초기화}는 ``탐험을 \textbf{초기값 하나로 예약}''하고,
    \kb{UCB}는 ``탐험을 \textbf{매 스텝의 불확실성으로 재계산}''
    한다.
  \end{block}
  \begin{itemize}
    \item 고정 환경: 둘 다 ``최선 팔 찾기''에 성공.
    \item \textbf{환경이 바뀌는 문제에서는 UCB의 ``계속
      재계산''만이 살아남음}.
    \item 두 전략 모두 ε-greedy의 ``정보와 무관한 고정 탐험''보다
      나은 방향 --- \textbf{불확실성에 비례하는 탐험}.
  \end{itemize}
  \vskip0.5em
  \textbf{2.2절 핵심 질문}: 같은 문제로 세 전략을 붙여 싸우면,
  정보와 무관하게 탐험하는 ε-greedy가 가장 \textbf{낭비}하고,
  팔의 순위가 뒤바뀌는 비정상 환경에서는 매 스텝 보너스를
  재계산하는 \textbf{UCB만이} 새 최선을 다시 찾아 적응.
\end{frame}
